\scanpage{6}
% L01–03
(on identifie $\mathbb E^*\times\mathbb E^{n-1}$ à l'ouvert de $\mathbb E^{\unsure{n}}$ dont les \deleted{\unsure{élé}} $k$-pts sont les $(a_0,\text{--},a_{n-1})\in k^n$ tels que $a_0\in k^*$ \ldots.).

% L04–06
\uline{Question} $X_n$ est-il réduit (i.e. lisse i.e. $\simeq\mathbb E^n$) ? Ou au contraire l'ens de ses pts réduits est il réduit à $X_n^{\mathrm{ét}}\simeq\mathbb E^*\times\mathbb E^{n-1}$ ?

% L07
Cas $n=1$, $X_1$ est-il réduit \inserted{i.e. lisse} au pt $0$ ?

% L08–13
Il revient au $=$ de dire que l'espace tangent de Zariski de $X_1$ en $0$ est de dim $1$ (\unsure{à présent|à priori} c'est de dim $\geq1$). \deleted{\unsure{Donc}} Il faut étudier les pts de $X_0$ à val. dans $k=\mathbb F_p[\unsure{t|T|J}]/t^2$ se réduisant suivant $0$, i.e. \unsure{c} les sous-groupes plats \inserted{$H$} de $G_{a\,k}$ tels que $H_0=\alpha_p$.

% L14–15
\diagram{\begin{tikzcd}[ampersand replacement=\&,column sep=small,row sep=small]
 H\overset{?}{\subset}\underset{\substack{\shortparallel\\G_a}}{G} \& H_0\subset G_0\arrow[r] \& \deleted{\unsure{G_0/H}}K_0\simeq G_0/H_0\\
 \& \overline H_0\subset G \& \overline K_0\simeq G/\overline H_0\simeq\overline G_0/\overline H_0
\end{tikzcd}}
% L16–20
$K=G/H$ est une \unsure{prol.|part.} de $K_0=G_0/H_0$ défini par une extension de $K_0$ par $\unsure{\mathfrak g K_0|\mathfrak L K_0|\mathbb V t_{K_0}|\mathscr L K_0|J K_0}$ (\unsure{don|dont} les splittings correspondent aux hom de $\overline K_0$ \inserted{$=G/\overline H_0$} avec $\deleted{\unsure{\overline K_0}}K\simeq G/H\deleted{\unsure{{}\simeq\overline G_0/\overline H_0}}$ \deleted{\unsure{)}} qui prolongent l'hom identique)

% L18–23
% slanted
\margin{Dégager \unsure{lemmes} pour calcul \unsure{des ces|de ces} espaces \unsure{tangents} de Zariski et obstructions à prolonger \unsure{ss-groupe} plat d'un \unsure{groupe} comm\ldots}
% L21–24
Image inverse de cette extension par $G_0\mathbin{\text{---}}K_0$ est l'obstruction à prolonger $G_0\overset{\unsure{\xi|\varepsilon}}{\mathbin{\text{---}}}K_0$ en $G\longrightarrow\deleted{\unsure{\overline K_0}}K$, or celle-ci se prolonge de façon canonique, \deleted{\unsure{en $G$}} donc cette extension
% archive sheet number: 5
